\newtcolorbox{BookSideNote}{
  enhanced,colback=white,colframe=black,boxrule=.6bp,arc=3bp,
  boxsep=0pt,left=7bp,right=7bp,top=5bp,bottom=5bp,
  before skip=0pt,after skip=0pt,drop fuzzy shadow,
  fontupper=\kaishu\small
}
\newcommand{\BookSideNoteClearance}{2bp}

\newsavebox{\BookNoteBody}
\newsavebox{\BookNoteAside}
\newdimen\BookNoteLastDepth
\newdimen\BookNoteResumeDepth

% Enter on the first body baseline. On exit the lower column determines
% the last line: either the actual body line, or a normal line whose
% lower edge is at the note bottom. Never add the height difference to
% prevdepth: that would eat the following baseline glue.
\newcommand{\BookSideNotePair}[3]{%
  \par\begingroup
  \sbox{\BookNoteBody}{%
    \begin{minipage}[t]{.52\linewidth}
      \ignorespaces#2\par
      \global\BookNoteLastDepth=\prevdepth
    \end{minipage}%
  }%
  \sbox{\BookNoteAside}{%
    \begin{minipage}[t]{.44\linewidth}
      \vspace{0pt}\begin{BookSideNote}#3\end{BookSideNote}%
    \end{minipage}%
  }%
  \setbox\BookNoteAside=\hbox{\raise\ht\BookNoteBody\box\BookNoteAside}%
  \BookNoteResumeDepth=\BookNoteLastDepth
  \ifdim\BookNoteResumeDepth<-999pt
    \BookNoteResumeDepth=\dp\strutbox
  \fi
  \ifdim\dp\BookNoteAside>\dp\BookNoteBody
    \dp\BookNoteAside=\dimexpr\dp\BookNoteAside+\BookSideNoteClearance\relax
    \BookNoteResumeDepth=\dp\strutbox
  \fi
  \noindent
  \if l#1\relax
    \box\BookNoteAside\hfill\box\BookNoteBody
  \else
    \box\BookNoteBody\hfill\box\BookNoteAside
  \fi
  \par\prevdepth=\BookNoteResumeDepth
  \endgroup
}

\newcommand{\BookPairedNote}[2]{\BookSideNotePair{r}{#1}{#2}}
\newcommand{\BookPairedNoteLeft}[2]{\BookSideNotePair{l}{#2}{#1}}

% In continuous prose, the same note bottom sets the end of the narrow
% measure. Include its shadow in the height that wrapfig measures.
\newenvironment{BookFlowNote}[2][r]{%
  \par
  \setlength{\intextsep}{0pt}%
  \setlength{\columnsep}{.04\linewidth}%
  \begin{wrapfigure}{#1}{.44\linewidth}%
    \begin{BookSideNote}#2\end{BookSideNote}%
    \vspace{\BookSideNoteClearance}%
  \end{wrapfigure}%
}{\par}

% Keep the note next to the introduction even inside the theorem list.
\newcommand{\ExampleSideNote}[2]{%
  \leavevmode\strut\par\smallskip
  \BookPairedNote{#1}{#2}%
}
